\sectionkicker{Source-backed technical notes}
\subsection*{単発LaunchからModel Family Lifecycleへ — GPT-5→5.1→5.2: Technical Notes}
本欄は記事本文で圧縮した一次資料上の情報を、比較・再検証しやすい形で整理したものである。時系列、確認済みの主張、留意点、一次資料URLを掲載する。
\medskip
\subsection*{Theme at a glance}
\addcontentsline{toc}{subsection}{Theme at a glance}
\begin{center}
\footnotesize
\begin{tabularx}{\linewidth}{@{}p{0.20\linewidth}p{0.10\linewidth}p{0.14\linewidth}X@{}}
\toprule
資料 & 位置づけ & 種別 & 時系列 \\
\midrule
GPT-5 release / developer API / safe-completions & 主要資料 & モデル & 2025-08-07 (モデル公開); 2025-08-07 (API公開); 2025-08-07 (SAFETY METHOD（公開）)  \\
GPT-5.1 & 補足資料 & モデル更新 & 2025-11-12 (モデル公開)  \\
GPT-5.1-Codex-Max & 補足資料 & モデル更新 & 2025-11-19 (CODING モデル公開（技術イベント）)  \\
GPT-5.2 & 補足資料 & モデル & 2025-12-11 (モデル公開)  \\
GPT-5.2-Codex & 補足資料 & モデル更新 & 2025-12-18 (CODING モデル公開（技術イベント）)  \\
\bottomrule
\end{tabularx}
\end{center}
\normalsize

\begin{technicalnote}{GPT-5 release / developer API / safe-completions}{主要資料}
\begingroup
% reader-facing Technical Notes break policy
\widowpenalty=10000
\clubpenalty=10000
\displaywidowpenalty=10000
\begin{tabularx}{\linewidth}{@{}>{\bfseries}p{0.22\linewidth}X@{}}
組織 & OpenAI \\
種別 & モデル
時系列 & 2025-08-07 (モデル公開); 2025-08-07 (API公開); 2025-08-07 (SAFETY METHOD（公開）) \\
\end{tabularx}
\smallskip
{\bfseries 一次資料から整理したtechnical points}
\begin{itemize}[leftmargin=1.5em,itemsep=0.35em]
\item \textbf{一次情報で確認できる事実}: 「GPT-5 release / developer API / safe-completions」について、一次資料で確認できる公開・提供・機能・時系列上の事実を要約した項目。
\item \textbf{一次情報で確認できる事実}: 「GPT-5 release / developer API / safe-completions」について、一次資料で確認できる公開・提供・機能・時系列上の事実を要約した項目。
\end{itemize}
\begin{minipage}{\linewidth}
% reader-facing Technical Notes generic boundary/source tail group
{\bfseries 読む際の境界}
\begin{itemize}[leftmargin=1.5em,itemsep=0.35em]
\item \textbf{分析上の留意点}: 「GPT-5 release / developer API / safe-completions」の扱いでは、一次資料で確認できる範囲、提供時点、評価条件を維持し、記載範囲を超えて一般化しない。
\end{itemize}
\begin{samepage}
{\bfseries 一次資料}
\begingroup\sloppy
\Urlmuskip=0mu plus 2mu\relax
\begin{itemize}[leftmargin=1.5em,itemsep=0.25em]
\item {\scriptsize\url{https://openai.com/index/gpt-5-safe-completions}}
\item {\scriptsize\url{https://openai.com/index/introducing-gpt-5}}
\item {\scriptsize\url{https://openai.com/index/introducing-gpt-5-for-developers}}
\end{itemize}
\endgroup
\end{samepage}
\end{minipage}
\endgroup
\end{technicalnote}
\medskip

\begin{technicalnote}{GPT-5.1}{補足資料}
\begingroup
% reader-facing Technical Notes break policy
\widowpenalty=10000
\clubpenalty=10000
\displaywidowpenalty=10000
\begin{tabularx}{\linewidth}{@{}>{\bfseries}p{0.22\linewidth}X@{}}
組織 & OpenAI \\
種別 & モデル更新
時系列 & 2025-11-12 (モデル公開) \\
\end{tabularx}
\smallskip
{\bfseries 一次資料から整理したtechnical points}
\begin{itemize}[leftmargin=1.5em,itemsep=0.35em]
\item \textbf{一次情報で確認できる事実}: OpenAIはGPT-5.1を発表し、ChatGPTの有料ユーザー向けに提供を開始した。
\end{itemize}
\begin{minipage}{\linewidth}
% reader-facing Technical Notes generic boundary/source tail group
{\bfseries 読む際の境界}
\begin{itemize}[leftmargin=1.5em,itemsep=0.35em]
\item \textbf{分析上の留意点}: ここで確認できるのは一次資料に記録された範囲の事実であり、ベンダー・プロジェクト・著者による評価を独立再現された結果として扱うものではない。
\end{itemize}
\begin{samepage}
{\bfseries 一次資料}
\begingroup\sloppy
\Urlmuskip=0mu plus 2mu\relax
\begin{itemize}[leftmargin=1.5em,itemsep=0.25em]
\item {\scriptsize\url{https://openai.com/index/gpt-5-1}}
\end{itemize}
\endgroup
\end{samepage}
\end{minipage}
\endgroup
\end{technicalnote}
\medskip

\begin{technicalnote}{GPT-5.1-Codex-Max}{補足資料}
\begingroup
% reader-facing Technical Notes break policy
\widowpenalty=10000
\clubpenalty=10000
\displaywidowpenalty=10000
\begin{tabularx}{\linewidth}{@{}>{\bfseries}p{0.22\linewidth}X@{}}
組織 & OpenAI \\
種別 & モデル更新
時系列 & 2025-11-19 (CODING モデル公開（技術イベント）) \\
\end{tabularx}
\smallskip
{\bfseries 一次資料から整理したtechnical points}
\begin{itemize}[leftmargin=1.5em,itemsep=0.35em]
\item \textbf{一次情報で確認できる事実}: OpenAIはCodex向けにGPT-5.1-Codex-Maxを公開し、長時間にわたるプロジェクト規模のエージェント型コーディング作業を想定したモデルと説明した。
\end{itemize}
\begin{minipage}{\linewidth}
% reader-facing Technical Notes generic boundary/source tail group
{\bfseries 読む際の境界}
\begin{itemize}[leftmargin=1.5em,itemsep=0.35em]
\item \textbf{分析上の留意点}: ここで確認できるのは一次資料に記録された範囲の事実であり、ベンダー・プロジェクト・著者による評価を独立再現された結果として扱うものではない。
\end{itemize}
\begin{samepage}
{\bfseries 一次資料}
\begingroup\sloppy
\Urlmuskip=0mu plus 2mu\relax
\begin{itemize}[leftmargin=1.5em,itemsep=0.25em]
\item {\scriptsize\url{https://openai.com/index/gpt-5-1-codex-max}}
\end{itemize}
\endgroup
\end{samepage}
\end{minipage}
\endgroup
\end{technicalnote}
\medskip

\begin{technicalnote}{GPT-5.2}{補足資料}
\begingroup
% reader-facing Technical Notes break policy
\widowpenalty=10000
\clubpenalty=10000
\displaywidowpenalty=10000
\begin{tabularx}{\linewidth}{@{}>{\bfseries}p{0.22\linewidth}X@{}}
組織 & OpenAI \\
種別 & モデル
時系列 & 2025-12-11 (モデル公開) \\
\end{tabularx}
\smallskip
{\bfseries 一次資料から整理したtechnical points}
\begin{itemize}[leftmargin=1.5em,itemsep=0.35em]
\item \textbf{一次情報で確認できる事実}: OpenAIはGPT-5.2をChatGPTとAPI向けに公開し、推論、長いコンテキスト、コーディング、画像理解を、専門業務やエージェント型ワークフロー向けの機能として位置付けた。
\end{itemize}
\begin{minipage}{\linewidth}
% reader-facing Technical Notes generic boundary/source tail group
{\bfseries 読む際の境界}
\begin{itemize}[leftmargin=1.5em,itemsep=0.35em]
\item \textbf{分析上の留意点}: ここで確認できるのは一次資料に記録された範囲の事実であり、ベンダー・プロジェクト・著者による評価を独立再現された結果として扱うものではない。
\end{itemize}
\begin{samepage}
{\bfseries 一次資料}
\begingroup\sloppy
\Urlmuskip=0mu plus 2mu\relax
\begin{itemize}[leftmargin=1.5em,itemsep=0.25em]
\item {\scriptsize\url{https://openai.com/index/introducing-gpt-5-2}}
\end{itemize}
\endgroup
\end{samepage}
\end{minipage}
\endgroup
\end{technicalnote}
\medskip

\begin{technicalnote}{GPT-5.2-Codex}{補足資料}
\begingroup
% reader-facing Technical Notes break policy
\widowpenalty=10000
\clubpenalty=10000
\displaywidowpenalty=10000
\begin{tabularx}{\linewidth}{@{}>{\bfseries}p{0.22\linewidth}X@{}}
組織 & OpenAI \\
種別 & モデル更新
時系列 & 2025-12-18 (CODING モデル公開（技術イベント）) \\
\end{tabularx}
\smallskip
{\bfseries 一次資料から整理したtechnical points}
\begin{itemize}[leftmargin=1.5em,itemsep=0.35em]
\item \textbf{一次情報で確認できる事実}: OpenAIはGPT-5.2-Codexを、長期的な推論と大規模なコード変換を想定したコーディングモデルとして公開した。
\end{itemize}
\begin{minipage}{\linewidth}
% reader-facing Technical Notes generic boundary/source tail group
{\bfseries 読む際の境界}
\begin{itemize}[leftmargin=1.5em,itemsep=0.35em]
\item \textbf{分析上の留意点}: ここで確認できるのは一次資料に記録された範囲の事実であり、ベンダー・プロジェクト・著者による評価を独立再現された結果として扱うものではない。
\end{itemize}
\begin{samepage}
{\bfseries 一次資料}
\begingroup\sloppy
\Urlmuskip=0mu plus 2mu\relax
\begin{itemize}[leftmargin=1.5em,itemsep=0.25em]
\item {\scriptsize\url{https://openai.com/index/introducing-gpt-5-2-codex}}
\end{itemize}
\endgroup
\end{samepage}
\end{minipage}
\endgroup
\end{technicalnote}
\medskip
